% Original exposition of the prime-pattern argument, prepared for Version 5.
% Requires amsmath, amssymb, amsthm, hyperref and the manuscript theorem environments.
% Publication-status qualifications in this file are intentional.
\section{Simultaneous primes with coprime adjacent norms}
\label{sec:v5-sun-proof}

This appendix gives the prime-pattern argument needed for the construction.
Its strategy is due to Sun~\cite{v5sun}. We supply the local calculations,
the two ranges of the sieve, the normalization of the ray-class densities,
and the order in which the parameters are chosen. One input has a different
publication status from the others. The Artin-symbol linear-forms theorem
of Frei and Tinkov\'a~\cite{v5ft}, Theorem~1.4 of arXiv version~3, is
unconditional and is listed by an author as accepted by \emph{Mathematika}.
We have not verified a final published version. We use that exact theorem
as an external input, and do not claim that the following proof has only
published external inputs. In particular, the proof below does not replace
their ideal-norm nilsequence argument.

Write $\mathrm N(z)=z\bar z$ for $z\in\mathbb Z[i]$, and put
$d_-(z)=\mathrm N(z-1)$ and $d_+(z)=\mathrm N(z+1)$.
An Artin or Frobenius element in this appendix is arithmetic Frobenius:
on residue fields it sends $x$ to $x^{\mathrm N\mathfrak p}$.

\begin{theorem}[The simultaneous prime input]\label{thm:v5-simultaneous-primes}
Let $c_1,\ldots,c_k$ be distinct integers, where $k\geq2$.
Let $q\geq1$, let $M$ be an odd squarefree positive integer with $7\nmid M$,
and let $H,Y>0$. There are positive integers $u,w$ such that
\[
 q\mid w,\qquad u/w^2>H,\qquad p_j=u+c_jw^2>Y
 \quad(1\leq j\leq k),
\]
where every $p_j$ is prime. Each $p_j$ has a Gaussian prime generator
$\pi_j$ of norm $p_j$ such that
\[
 \pi_j\equiv9\pmod {16},\qquad
 \pi_j\equiv1+3i\pmod7,\qquad
 \pi_j\equiv i\pmod M.
\]
No odd prime divides both $d_-(\pi_j)$ and $d_+(\pi_j)$.
For $j\ne h$, no odd prime divides both $d_\epsilon(\pi_j)$ and
$d_\eta(\pi_h)$, for any $\epsilon,\eta\in\{-,+\}$.
Moreover,
\[
 v_2(d_-(\pi_j))=6,\quad v_2(d_+(\pi_j))=2,\quad
 \gcd(d_-(\pi_j),d_+(\pi_h))=4\quad\hbox{for all }j,h,
\]
and all these norms are coprime to $7M$.
\end{theorem}

The quantifiers concern fixed data. No estimate uniform in $q,M,k$ or the
$c_j$ is asserted. We first prove the result for $q=1$. Replacing the
coefficients by $q^2c_j$ and then replacing the parameter $v$ by $w=qv$
proves the general case, since the ratio before this substitution can be
made greater than $q^2H$.

\subsection{Local counts and the medium primes}\label{subsec:v5-local-sieve}

Fix a nonzero integer $\kappa$ and write
\[
 Q_\kappa=a^2+b^2-c^2-d^2-\kappa w^2,\quad
 G_\epsilon=(a+\epsilon)^2+b^2,\quad
 H_\eta=(c+\eta)^2+d^2,
 \qquad \epsilon,\eta\in\{1,-1\}.
\]
We use $\chi(\ell)=(-1/\ell)$ for odd primes $\ell$.

\begin{lemma}[Finite-field counts]\label{lem:v5-local-counts}
If $\ell\nmid2\kappa$, then $Q_\kappa=0$ has $\ell^4-1$ nonzero
solutions over $\mathbb F_\ell$. At most $4\ell^2$ of its solutions satisfy
$G_\epsilon=H_\eta=0$. If $\chi(\ell)=-1$, the latter set has exactly one
element.

Let $\gamma_1,\ldots,\gamma_k$ be distinct integers, including $0$ and
$-\kappa$, and put $L_j=a^2+b^2+\gamma_jw^2$. Outside a fixed set of primes
containing the divisors of $2\kappa\prod_{j<h}(\gamma_j-\gamma_h)$, the
proportion $g(\ell)$ of nonzero zeros of $Q_\kappa$ on which some $L_j$
vanishes is
\begin{equation}\label{eq:v5-local-g}
 g(\ell)=
 \frac{(\ell-1)(\ell-\chi(\ell))
 [2\nu_0+(k-2)(\ell-\chi(\ell))]+\nu_0^2-1}{\ell^4-1},
 \qquad \nu_0=\ell+(\ell-1)\chi(\ell).
\end{equation}
Consequently $g(\ell)=(k+2\chi(\ell))/\ell+O_k(\ell^{-2})$.
\end{lemma}
\begin{proof}
The norm $x^2+y^2$ has $\nu_0$ representations of zero and
$\ell-\chi(\ell)$ representations of every nonzero element. This follows
either by factoring over $\mathbb F_\ell$ when $\chi=1$, or by the norm map
from $\mathbb F_{\ell^2}$ when $\chi=-1$. Summing these fibres for
$Q_\kappa$ gives $\ell^4$ solutions including the origin.
In the split case, each of the equations $G_\epsilon=0$ and $H_\eta=0$
is a union of two lines. On each of the four pairs of lines the remaining
equation is
$-2\epsilon a+2\eta c-\kappa w^2=0$, leaving at most $\ell^2$ choices.
In the inert case these equations force
$(a,b,c,d)=(-\epsilon,0,-\eta,0)$, and then $w=0$.

For $w\ne0$, the conditions $L_j=0$ are disjoint. For the two distinguished
coefficients one norm is zero and the other is nonzero, giving
$2\nu_0(\ell-\chi)$ choices. Each other coefficient gives
$(\ell-\chi)^2$ choices. There are $\ell-1$ possible nonzero $w$.
For $w=0$, both norms are zero, giving $\nu_0^2-1$ nonzero vectors.
This proves \eqref{eq:v5-local-g}. In particular, when $k=2$ and
$\chi=-1$, its value is $2/(\ell^2+1)$, not zero.
\end{proof}

\begin{lemma}[Primitive congruence count]\label{lem:v5-primitive-count}
Let $\psi$ be a smooth nonnegative compactly supported function away from
the origin, and let $m$ be squarefree and prime to $2\kappa$.
For each $\ell\mid m$ prescribe an arbitrary subset $\Omega_\ell$ of the
nonzero zeros of $Q_\kappa$ over $\mathbb F_\ell$. The weighted number of
primitive integral vectors on $Q_\kappa=0$ with these reductions equals
\begin{equation}\label{eq:v5-primitive-count}
 C_\kappa\sigma_\infty(\psi)B^3
 \prod_{\ell\mid m}\frac{|\Omega_\ell|}{\ell^4-1}
 +O_{\kappa,\psi,\varepsilon}
       (B^{5/2+\varepsilon}m^{5/2+\varepsilon}).
\end{equation}
The constant $C_\kappa$ is positive and independent of $m$ and the subsets.
\end{lemma}
\begin{proof}
Apply Browning--Loughran~\cite{v5bl}, Theorem~4.1, to $Q_\kappa$ in five
variables. It is isotropic, since $(1,0,1,0,0)$ is a zero, and nonsingular,
since $\kappa\ne0$. Its discriminant has no prime divisor outside
$2\kappa$. The theorem permits $n=5$ and arbitrary nonzero residue subsets;
its error is uniform in the modulus and in these subsets.
Impose primitivity by $\sum_{h\mid(a,b,c,d,w)}\mu(h)$.
If $(h,m)>1$ there are no permitted vectors. Otherwise the subsets become
$h^{-1}\Omega_\ell$, of the same cardinalities, and the scale becomes $B/h$.
The sum of the error terms is convergent. Extending the main-term sum to
infinity costs $O(B)$.
At a good prime the unrestricted local density is
\[
 \sigma_\ell=\frac{1-\ell^{-4}}{1-\ell^{-3}}.
\]
Indeed, nonsingular primitive zeros lift with four free coordinates, and
the vectors divisible by $\ell$ contribute the geometric factor
$\ell^{-3}$. The primitive density is therefore $1-\ell^{-4}$.
This turns the factor $|\Omega_\ell|/\ell^4$ of the unprimitive count into
$|\Omega_\ell|/(\ell^4-1)$ in \eqref{eq:v5-primitive-count}.
The remaining finitely many local factors are positive because the form
has a nonsingular rational zero. Empty subsets give the assertion directly.
\end{proof}

In using this count, the two norms lie between fixed positive multiples of
$B^2$, and $|w|\leq\tau B$, where $0<\tau\leq\tau_0$.
A smooth majorant can be chosen with
$\sigma_\infty(\psi)\ll\tau$ uniformly in $\tau$: integrate first in the
two radial norm variables, on which the derivative of their difference is
bounded away from zero. The constants in error terms may depend on
derivatives of this majorant, and hence on $\tau$.

\begin{lemma}[Medium-prime bound]\label{lem:v5-medium-sieve}
For fixed coefficients and sufficiently large fixed $L$, the contribution
of prime patterns with a common adjacent-norm divisor
$L<\ell\leq B^{1/100}$ is at most
\begin{equation}\label{eq:v5-medium-sieve}
 \frac{C\tau B^3}{(\log B)^k L\log L}+O_\tau(B^{27/10}).
\end{equation}
The count is unweighted by logarithms, and $C$ is independent of $L,\tau$.
The bound covers both distinct positions and the two signs at one position.
\end{lemma}
\begin{proof}
For two positions, use $Q_\kappa$ with their fixed nonzero coefficient
difference and impose $G_\epsilon=H_\eta=0\pmod\ell$.
For squarefree $d$ supported on good primes other than $\ell$, impose also
that some $L_j$ vanishes at every prime dividing $d$. The preceding lemmas
give a sieve sequence
\[
 A_{\ell,d}=X_\ell g(d)+r_{\ell,d},\qquad
 X_\ell\ll\tau B^3/\ell^2,\qquad
 |r_{\ell,d}|\ll_{\tau,\varepsilon}
 B^{5/2+\varepsilon}(\ell d)^{5/2+\varepsilon}.
\]
Enlarge the fixed exceptional set so that $0<g(p)<1$ outside it.
Set $h(p)=g(p)/(1-g(p))$ and extend $h$ multiplicatively to squarefree
integers. The elementary Selberg upper sieve has denominator
\[
 \mathcal G_\ell(z)=\sum_{\substack{d<z\ \text{squarefree}\\
                         (d,\ell\prod_{p\in S}p)=1}}h(d)
 \gg (\log z)^k.
\]
Here is the dimension calculation. Williams~\cite{v5williams}, Theorem~1
with fixed modulus $4$ and residues $1,3$, gives Mertens' estimates in the
two reduced classes. Taking logarithms gives
$\sum_{p\leq y}g(p)=k\log\log y+O(1)$ and
$\sum_{p\leq y}g(p)\log p\ll\log y$.
For the latter bound only the elementary Chebyshev estimate is needed:
$\prod_{n<p\leq2n}p\mid\binom{2n}{n}<4^n$ gives
$\vartheta(x)\ll x$ by dyadic summation, and partial summation gives
$\sum_{p\leq y}(\log p)/p\ll\log y$.
In the Euler product $\prod_{p\leq y}(1+h(p))$, normalize the weights
$h(d)$ to a probability distribution. Its expected value of $\log d$ is
$\sum_{p\leq y}g(p)\log p$. Taking $y$ to be a sufficiently small fixed
power of $z$ leaves at least half of the mass on $d<z$, by Markov's
inequality. Removing one prime $\ell$ changes the estimate by a bounded
factor.

For completeness the sieve weights may be taken as
\[
 \lambda_d=\frac{\mu(d)}{\mathcal G_\ell(z)}
 \prod_{p\mid d}(1-g(p))^{-1}
 \sum_{\substack{r<z/d\ (r,d)=1\ r\ \text{squarefree}\\
                      (r,\ell\prod_{p\in S}p)=1}}h(r).
\]
They have $\lambda_1=1$, $|\lambda_d|\ll_\varepsilon d^\varepsilon$,
and direct expansion gives
$\sum_{d,e}\lambda_d\lambda_e g([d,e])=1/\mathcal G_\ell(z)$.
Thus the number with no prime divisor below $z$ in any $L_j$ is at most
\[
 \frac{C\tau B^3}{\ell^2(\log z)^k}
 +O_{\tau,\varepsilon}
  (B^{5/2+\varepsilon}\ell^{5/2+\varepsilon}z^{7+4\varepsilon}).
\]
All the prime values under consideration exceed $z$. Put $z=B^{1/100}$
and sum over $L<\ell\leq B^{1/100}$. The main term uses
$\sum_{\ell>L}\ell^{-2}\ll1/(L\log L)$, another consequence of the
same Chebyshev bound; for example
$\varepsilon=1/1000$ makes the error smaller than $B^{27/10}$.

At one position, an odd common divisor of $(a-1)^2+b^2$ and
$(a+1)^2+b^2$ forces $a=0$ and $b^2=-1\pmod\ell$.
There are at most $2\ell$ such triples $(a,b,w)$ over $\mathbb F_\ell$.
For $w\ne0$ the zero sets of the $k$ prime forms are disjoint, and the
ordinary three-dimensional local sieve density is
$(k+\chi(\ell))/\ell+O_k(\ell^{-2})$.
Counting residue classes in a smooth box gives the corresponding
three-dimensional congruence estimate. The same Selberg weights and the
same, weaker error in \eqref{eq:v5-medium-sieve} apply.
\end{proof}

\subsection{The large-prime tail}\label{subsec:v5-large-sieve}

We use two exact inputs from Browning--Heath-Brown~\cite{v5bhb}.
Their Lemma~2.1 says that coprime polynomials in $\mathbb Z[X_1,\ldots,X_m]$
of bounded degree and coefficient height at most $H$ have at most
\begin{equation}\label{eq:v5-box-sieve}
 O\left(\frac{V\log(VH)}{L\log L}
           +\frac{V\log(VH)}{B_{\min}}\right)
\end{equation}
integer points in a box of side parameters $B_i\geq2$ at which a prime
greater than $L\geq2$ divides every polynomial; $V=\prod B_i$.
``Coprime'' here excludes a common constant prime factor as well as a
common nonconstant factor. Their Lemma~4.1 covers the zeros modulo a
squarefree $q$ of a rank-$m$ quadratic form by at most
$(3m)^{\omega(q)}q^{m-2}$ lattices of determinant $q^{m-1}$.
The number with largest Euclidean successive minimum exceeding $T$ is
\begin{equation}\label{eq:v5-lattice-tail}
 O_m((3m)^{\omega(q)}q^{2m-3}T^{-m}).
\end{equation}
In this statement every prime dividing $q$ must be odd and preserve the
rank of the form. We include the word ``odd'' explicitly. Rank preservation
alone is not the good-prime convention of the cited lemma.

\begin{lemma}[Large common divisors on the quadric]\label{lem:v5-large-sieve}
For fixed $\kappa\ne0$ and $\epsilon,\eta\in\{1,-1\}$, the number of
integral vectors of size $O(B)$ satisfying $Q_\kappa=0$ and
$\ell\mid G_\epsilon,H_\eta$ for some prime $\ell>L$ is
\begin{equation}\label{eq:v5-large-sieve}
 O_{\kappa,\alpha}\left(
    \frac{B^{3+\alpha}}{L\log L}+B^{11/4+\alpha}\right)
 \qquad(\alpha>0).
\end{equation}
\end{lemma}
\begin{proof}
Sign changes of $a,c$ reduce to $\epsilon=\eta=1$.
Put $x_0=a+c$, $x_1=a-c$, $x=b-d$, $z=b+d$, and
$R=\kappa w^2-xz$. Then $x_0x_1=R$.
Parity restrictions can be discarded in an upper bound.

First suppose that $\ell\mid x_0,x_1$. Then $\ell^2\mid R$.
For each fixed integer $v=O(B^2)$, the equation $R=v$ has
$O_{\kappa,\alpha}(B^{1+\alpha})$ triples: for each $w$, count the
factorizations $xz=\kappa w^2-v$ by the divisor bound; the zero right-hand
side occurs for at most two $w$ and contributes $O(B)$ each.
If $R\ne0$, it has $O_\alpha(B^\alpha)$ factorizations $x_0x_1=R$.
Summing over $R=h\ell^2\ne0$ gives
$O(B^{3+2\alpha}/(L\log L))$.
The case $R=0$, including $x_0x_1=0$, contributes $O(B^{2+\alpha})$.

In the remaining cases choose $t=x_1$ or $t=x_0$ such that $\ell\nmid2t$.
Take $t=x_1$ first. Multiplication by the invertible denominator $4t^2$
turns the two norm congruences into divisibility of
\[
 P=\kappa w^2+2t,\qquad
 K=(\kappa w^2-xz+t^2+2t)^2+t^2(x+z)^2.
\]
The identity
\[
 K-(t^2+x^2)(t^2+z^2)=P(P+2t^2-2xz)
\]
shows that it suffices to sieve the coprime pair
$P,(t^2+x^2)(t^2+z^2)$. They have no common nonconstant factor over
$\mathbb Q$, because the first depends nontrivially only on $w$, whereas
the second is independent of $w$. Their contents have no common prime,
since the second has coefficient $1$ on $x^2z^2$.
For $t=x_0$ the pair is
\[
 P'=\kappa(t+2)w^2-2xz,\qquad K.
\]
If $t\ne-2$, $P'$ is an irreducible ternary quadratic form of rank three
and $K$ has nonzero constant term $t^2(t+2)^2$; homogeneity rules out
divisibility. If $t=-2$, neither $x$ nor $z$ divides $K$.
The coefficient $-2$ of $xz$ excludes a common odd content prime.

In either chart, $t\mid R$. Put
\[
 S_R=\{p:p\mid2\kappa\},\qquad
 q=\prod_{\substack{p\mid t\ p\notin S_R}}p.
\]
The form $R$ has rank three at every prime in this product, and those primes
are odd. For a fixed $q$, the number of eligible $t=O(B)$ is
$O_{\kappa,\alpha}(B^\alpha)$ after an arbitrarily small adjustment of
$\alpha$. One way to see this is to majorize the count of integers
supported on $S_R\cup\{p:p\mid q\}$ by
$B^\alpha\prod_{p\in S_R\ \text{or}\ p\mid q}(1-p^{-\alpha})^{-1}$,
and absorb its $q^{o(1)}$ factor.

Apply \eqref{eq:v5-lattice-tail} with $m=3$ to $q\mid R$.
Each covering lattice contains $q\mathbb Z^3$, so all its successive minima
are $O(q)=O(B)$. It therefore has $O(B^3/q^2)$ points in the box.
The lattices with largest minimum greater than $T$ contribute, after the
sums over $q$ and $t$,
\[
 O_{\kappa,\alpha}(B^{5+2\alpha}T^{-3}).
\]
For the other lattices choose a basis with lengths comparable to the
successive minima. The resulting rectangular box has volume $O(B^3/q^2)$
and minimum side parameter $\gg B/T$.
The substitution into the polynomial pair is invertible over $\mathbb Q$.
For $\ell\nmid2t$ its determinant, of absolute value $q^2$, is invertible
modulo $\ell$ too. Thus no new common content prime outside $2t$ appears.
Remove common content supported on $2t$, which cannot affect the primes
being counted. Coefficient heights are powers of $B$ with fixed exponents.
Formula \eqref{eq:v5-box-sieve} bounds this lattice's contribution by
\[
 O\left(\frac{B^3\log B}{q^2L\log L}
                 +\frac{B^2T\log B}{q^2}\right).
\]
There are at most $9^{\omega(q)}q$ lattices. Summation gives
$O(B^{3+3\alpha}/(L\log L)+B^{2+3\alpha}T)$.
Choose $T=B^{3/4}$ and rename the small exponent to obtain
\eqref{eq:v5-large-sieve}.
\end{proof}

\begin{proposition}[Weighted obstruction tail]\label{prop:v5-obstruction-tail}
For fixed distinct $c_j$ there are $\tau_0,C>0$ such that, for every fixed
$0<\tau\leq\tau_0$, the logarithmically weighted number of pairs
\[
 N/2\leq u\leq N,\qquad 1\leq w\leq\tau\sqrt N,
 \qquad p_j=u+c_jw^2\text{ all prime},
\]
having any of the odd common divisors forbidden in
Theorem~\ref{thm:v5-simultaneous-primes} greater than $L$, is at most
\begin{equation}\label{eq:v5-obstruction-tail}
 \frac{C\tau N^{3/2}}{L\log L}
       +O_\tau(N^{3/2-1/800}).
\end{equation}
The weight is $\prod_j\log p_j$, and the bound allows every Gaussian
generator, so it also holds for any prescribed choice of generators.
\end{proposition}
\begin{proof}
Take $B=\sqrt N$ and combine the preceding bounds at $L_0=B^{1/100}$.
The medium bound loses at most $(\log B)^k$ on inserting the weight.
The large bound with sufficiently small $\alpha$ is
$O(B^{3-1/200})$ after the logarithmic weight is inserted.
At one position use the polynomial pair $a,b^2+1$ in three variables
in \eqref{eq:v5-box-sieve}.
There are only finitely many pairs of positions and signs, and a split
rational prime has eight Gaussian generators. These factors are absorbed
in $C$. Every vector arising from prime norm is primitive, since a common
integer divisor greater than one would have its square divide that prime.
Weakening the power saving to $B^{3-1/400}$ gives the displayed estimate.
\end{proof}

\subsection{Fixed Artin conditions in quadratic patterns}
\label{subsec:v5-artin-transfer}

\begin{remark}[A norm exponent in the imported proof]
\label{rem:v5-ft-norm-repair}
The proof of Frei--Tinkov\'a's Lemma~9.3, reused in their Section~10,
chooses a majorant with the $D$-linear-forms condition for $D=4^u$ and
records $\|\nu-1\|_{U^{2u}}=o(1)$. Its next step applies the relative
inverse theorem to the $U^{u+1}[N]$ norm. Theorem~5.1 of
Dodos--Kanellopoulos~\cite{v5dk} requires
$\|\nu-1\|_{U^{2s}(\mathbb Z/N'\mathbb Z)}\leq\eta$ for the target
$U^s[N]$ norm. Thus the displayed choice does not supply that hypothesis.
Request instead $D=4^{u+1}$ in Frei--Tinkov\'a's Lemma~9.1. Expanding the
centered $(2u+2)$-cube uses at most $2^{2u+2}=4^{u+1}$ pairwise
nonproportional forms, in $2u+3\leq4^{u+1}$ variables for $u\geq1$,
with coefficients in $\{0,1\}$. The stronger linear-forms condition gives
$\|\nu-1\|_{U^{2u+2}}=o(1)$ and supplies the missing hypothesis.
All complexities remain fixed before $N$ tends to infinity.

For clarity, the inverse theorem takes a prime $N'\in[CN,2CN]$ with
$C\geq20$, a nonnegative majorant on $\mathbb Z/N'\mathbb Z$, and a real
function $f$ on $[N]$ with $|f|\leq\nu$ and $\|f\|_{U^s[N]}\geq\delta$.
Its constants depend on $s,C,\delta$ and its conclusion is correlation
with a bounded Lipschitz nilsequence on an $(s-1)$-step nilmanifold.
The truncation and positive rescaling already present in Lemma~9.3
put its real centered prime weight in this form, with $s=u+1$.
This change repairs that theorem application. It does not replace or
independently verify the ideal-norm estimates in their Theorem~2.4.
The exact inverse-theorem statement was checked in arXiv:1610.00487v4;
the work is published, but its final typeset statement was not compared.
\end{remark}

Here the field is fixed throughout the limiting argument. Let $K/\mathbb Q$
be finite Galois, put $G=\operatorname{Gal}(K/\mathbb Q)$, and let
$C_1,\ldots,C_k$ be conjugacy-invariant subsets of $G$.
Let $D$ be divisible by the conductor of $K^{\rm ab}$.
For a unit $b\pmod D$, define
\begin{equation}\label{eq:v5-mu}
 \mu_j(b)=\frac{\#\{g\in C_j:g|_{K^{\rm ab}}=\sigma_b|_{K^{\rm ab}}\}}
                  {[K:K^{\rm ab}]},
 \qquad \sigma_b(\zeta_D)=\zeta_D^b,
\end{equation}
and set $\mu_j(b)=0$ for nonunits. Define $\Lambda'_j(p)=\log p$ for
unramified primes with arithmetic Frobenius in $C_j$, and zero otherwise.

\begin{proposition}[Fixed-field polynomial count]\label{prop:v5-artin-patterns}
There is a function $\omega(N)\to0$, depending on the fixed data, such that
whenever $\omega(N)\sqrt N\leq T=o(\sqrt N)$,
\begin{equation}\label{eq:v5-artin-patterns}
 \sum_{u\leq N}\sum_{1\leq w\leq T}
       \prod_{j=1}^k\Lambda'_j(u+c_jw^2)
       =\mathfrak S\,NT+o(NT),
\end{equation}
where
\begin{equation}\label{eq:v5-artin-series}
 \mathfrak S=
 \left(\frac D{\varphi(D)}\right)^k
 \mathbb E_{u,w\bmod D}\prod_j\mu_j(u+c_jw^2)
 \prod_{p\nmid D}\beta_p,
 \quad
 \beta_p=\mathbb E_{u,w\bmod p}
       \prod_j\frac p{p-1}1_{p\nmid u+c_jw^2}.
\end{equation}
The same statement on $N/2\leq u\leq N$ has main term
$\mathfrak SNT/2$.
\end{proposition}
\begin{proof}
We specify how the linear Artin theorem and the published polynomial
transference results combine. Frei--Tinkov\'a~\cite{v5ft}, Theorem~1.4,
applies to any fixed finite-complexity system of affine-linear forms,
with bounded linear coefficients and bounded constant coefficients divided
by the ambient scale, uniformly over convex regions in a fixed box.
Its conductor factor is precisely \eqref{eq:v5-mu}; a finite union of
conjugacy classes follows by addition. This is their unconditional theorem,
not their primitive-root theorem, Theorem~1.2.

Choose a small positive $\rho$ and first keep $v$ fixed. Put
\[
 W=D\prod_{\substack{p\leq v\\p\nmid D}}p,\quad
 A=W^{1/\rho},\quad n=\lfloor N/W\rfloor,\quad R=n^\rho,
\]
and, for unit $b\pmod W$, put
\[
 f_{j,b}(x)=\frac{\varphi(W)}W\Lambda'_j(Wx+b)1_{R\leq x\leq n},
 \qquad g_{j,b}=f_{j,b}-\mu_j(b).
\]
View these functions on $\mathbb Z/n\mathbb Z$, with representatives
$1,\ldots,n$. If $a\geq2$ and $C_1<C_2<C_3$ are fixed, the linear Artin
theorem gives
\begin{equation}\label{eq:v5-global-gowers}
 \|g_{j,b}\|_{U^a_{q[J]}}^{2^a}
       \ll_a(\log v)^{-1}+o_v(1),
 \quad q\leq A^{C_1},\quad
 nA^{-C_3}\leq J\leq nA^{-C_2}.
\end{equation}
The steps in this norm are the difference multiset $q(h-h')$, with
$h,h'\in[J]$, as in Tao--Ziegler. On partitioning by the wrap-around
integers, its cube expansion consists of the affine forms
\[
 W\left(x+q\sum_{r=1}^a\nu_r(h_r-h'_r)-m_\nu n\right)+b,
 \qquad \nu\in\{0,1\}^a.
\]
They are pairwise nonproportional: the $x$ coefficient forces the ratio to
be one, and the $h_r$ coefficients then distinguish their vertices.
The wrap-around and interval conditions define finitely many convex
regions. For fixed $v$ every scale ratio is fixed, so the theorem's
absolute error is $o_v(nJ^{2a})$.
At $p\nmid Wq$ the local factor is $1+O_a(p^{-2})$; at $p\mid q$ it is
$1+O_a(p^{-1})$. Since $\log q=O(v)$, there are $O(v/\log v)$ prime
divisors of $q$ greater than $v$. The product of these factors is
$1+O_a(1/\log v)$. Centering cancels the common main term. Removing
$x<R$ costs $o_v(1)$, as do prime powers if the linear theorem is stated
with their von Mangoldt weights. This proves \eqref{eq:v5-global-gowers}.

We now use the arbitrary-function dense-model theorem of
Tao--Ziegler~\cite{v5tz}. Their published Theorem~11 and Corollary~3
give a single bounded model accurate in both global Gowers norms and the
local polynomial box norms. Their Theorem~9 is the bounded-function
concatenation statement, and their Theorem~13 bounds the polynomial
average by one of those local norms. All majorants here are the ordinary
prime majorants, since $0\leq\Lambda'_j\leq\Lambda'$.
The fixed factor $D$ in $W$ is allowed in their divisor-sum construction:
for sufficiently large $v$, $W=E\prod_{p\leq v}p$ with a fixed $E$ whose
prime factors are at most $v$. A unit residue makes every small-prime
divisibility condition identically false, including at the extra prime
powers. At the remaining primes the divisor-sum local factors are
unchanged. Polynomial coefficient bounds absorb the fixed factor $E$.
Thus the polynomial-forms estimates for the majorant have the same
$1+o(1)$ conclusion.

The following choice records the size hypotheses of these results.
First choose their local box order $d'$ and a coefficient exponent $C$;
choose an integer $B_0>C+2$, then a sufficiently large global order $a$,
and finally a sufficiently small $\rho$. Set $\rho_* =\rho/B_0$ and
$\mathcal A=A^{B_0}$. Put
$M_0=\lfloor T/W\rfloor$, $M_1=\lfloor M_0/A\rfloor$.
For $T/\sqrt N\geq4\sqrt W/A$ one has
\[
 A^{-1}\sqrt n\leq M_0\leq\sqrt n,
 \qquad A^{-3}\sqrt n\leq M_1\leq\sqrt n.
\]
The model at cutoff $n^{\rho_*}$ dominates the function already truncated
at $n^\rho$. Corollary~3 supplies approximation for the degree-one local
step polynomials, whose coefficients are $O(A^C)$, and simultaneously for
$q\leq A^{B_0a}$ with
$J=\lfloor A^{-2B_0a}M_1^2\rfloor$.
In particular
\[
 nA^{-2B_0a-7}\leq J\leq nA^{-2B_0a}.
\]
This is within \eqref{eq:v5-global-gowers}. The bounded model minus
$\mu_j(b)$ therefore has small global norms. Theorem~9 makes its local
norm small; its smallness threshold is independent of $\mathcal A$.
The model approximation transfers this conclusion to $g_{j,b}$, and
Theorem~13 makes every polynomial average containing a centered factor
$o(1)$. These are the reasons for choosing the complexities before $v$.

For the final $W$-trick write $u=Wx+b$ and $w=Wm+c$.
If $b_j$ is the residue of $b+c_jc^2$ modulo $W$, the resulting integer
polynomial is
\[
 c_jWm^2+2c_jcm+(b+c_jc^2-b_j)/W.
\]
Its leading coefficient differences are nonzero, and its coefficients
are $O(W)$. Expand the product into its means and centered factors and
apply the preceding conclusion. End intervals of length $O(T^2)=o(N)$
are harmless by the ordinary majorant. Summing residue classes does not
amplify the error: their normalized admissible mass is
$\prod_{p\mid W}\beta_p=O(1)$.
Indeed, outside the fixed coefficient-difference primes,
\[
 \beta_p=\left(\frac p{p-1}\right)^k
                \frac{(p-1)(p-k+1)}{p^2}=1+O_k(p^{-2}).
\]
This also proves convergence of the product in
\eqref{eq:v5-artin-series}.
Let $N$ tend to infinity for each fixed $v$ and then let $v$ tend to
infinity sufficiently slowly. A diagonal choice supplies $\omega(N)$
and yields \eqref{eq:v5-artin-patterns}. Subtract the corresponding
formula with upper limit $N/2$, enlarging $\omega$ if necessary, to
obtain the last assertion.
\end{proof}

\begin{remark}[Published numbering]\label{rem:v5-tz-numbering}
The labels in the published 2018 paper differ from those in arXiv
1603.07817v1. The relevant dictionary is
\[
\begin{array}{c|ccccc}
 \text{arXiv} &\text{Theorem 1.4}&\text{Theorem 2.5}&\text{Theorem 3.1}
               &\text{Corollary 3.9}&\text{Theorem 4.1}\\
 \text{published}&\text{Theorem 4}&\text{Theorem 9}&\text{Theorem 11}
               &\text{Corollary 3}&\text{Theorem 13}
\end{array}
\]
The published locations are respectively pages $4,21,28,38,41$.
The preprint's equations (3.5) and (4.1) are published equations (36)
and (57). The dense model used above is the theorem's assertion for an
arbitrary function bounded by the prime weight. The special model $1$
for the ordinary prime weight is not an assertion about Artin weights.
\end{remark}

\subsection{The ray-class lower bound}\label{subsec:v5-ray-density}

For ray class fields we use the published existence and reciprocity
statements in Kopp--Lagarias~\cite{v5kl}, Theorems~1.1, 1.3 and 7.2,
and the modulus exact sequence of Theorem~6.5, specialized to the maximal
order $\mathbb Z[i]$. There are no real places. This replaces an appeal
to unpublished course notes. In this principal ideal domain the ray group
for the finite modulus $H\mathbb Z[i]$ is
\begin{equation}\label{eq:v5-ray-group}
 (\mathbb Z[i]/H\mathbb Z[i])^\times/\{1,-1,i,-i\}.
\end{equation}
It also follows directly by sending a residue $r$ to the ideal class of
$(r)$: surjectivity is principality, and the kernel consists exactly of
unit residues. The Artin map sends this ideal class to arithmetic
Frobenius. Conjugation acts by $r\mapsto\bar r$.

Choose a fixed $L_1$ larger than all coefficient-difference primes and
larger than $11$. Let $H_0$ be divisible by $16\cdot7M$ and by all odd
primes at most $L_1$. Prescribe a unit residue $r\pmod{H_0}$ which is
$9$ modulo $16$, $1+3i$ modulo $7$, and $i$ modulo every other odd
prime dividing $H_0$. Compatible higher powers can be chosen arbitrarily.
The associated rational-prime Frobenius condition is the two-element
conjugacy class $C_0=\{h_r,h_{\bar r}\}$ in the Galois group $G_0$ of
the ray field $K_0$ over $\mathbb Q$.
The two elements are distinct: an equality of their ray classes would
require $r=\epsilon\bar r\pmod{H_0}$ for a Gaussian unit $\epsilon$;
modulo $16$ this forces $\epsilon=1$, contradicted modulo $7$.
The field is Galois over $\mathbb Q$, because conjugation preserves the
modulus and the defining ray subgroup.

Let $D_0$ be the conductor of $K_0^{\rm ab}$ and choose a unit
$b_0\pmod{D_0}$ inducing the restriction of $h_r$.
The value in \eqref{eq:v5-mu} is then
$\mu_0=|C_0|/|G_0'|>0$.
For $L\geq L_1$ let
$T_L=\prod_{L_1<\ell\leq L,\,\ell\nmid H_0}\ell$ and let $K_L$ be the
ray field for $H_0T_L\mathbb Z[i]$. For each new prime write
\[
 B_\ell=(\mathbb Z[i]/\ell\mathbb Z[i])^\times,
 \qquad N_\ell=\ker(B_\ell\xrightarrow{\mathrm N}\mathbb F_\ell^\times).
\]
The norm map is onto and $|N_\ell|=\ell-\chi(\ell)$.

\begin{lemma}[Abelian fibres and their surviving proportion]
\label{lem:v5-ray-fibres}
Put $N_L=\prod_{\ell\mid T_L}N_\ell$ and
$G_L=\operatorname{Gal}(K_L/\mathbb Q)$.
Then
\[
 N_L\subseteq G_L',\qquad
 G_L'/N_L\simeq G_0',\qquad
 K_L^{\rm ab}=K_0^{\rm ab}\mathbb Q(\zeta_{T_L}).
\]
Let $C_L$ consist of those elements above $C_0$ whose normalized Gaussian
residue $z$ satisfies $\mathrm N(z-1)\mathrm N(z+1)\not\equiv0\pmod\ell$
at every $\ell\mid T_L$. For the density $\mu_L$ of this set,
\begin{equation}\label{eq:v5-ray-density}
 \mu_L(b)\geq\mu_0
       \prod_{\ell\mid T_L}\left(1-\frac4{\ell-1}\right)
 \qquad (b\equiv b_0\pmod{D_0},\ (b,D_0T_L)=1).
\end{equation}
\end{lemma}
\begin{proof}
In the CRT presentation of \eqref{eq:v5-ray-group}, the unit group
embeds in the $H_0$ factor, because it already embeds modulo $16$.
Once the representative above $r$ or $\bar r$ is normalized, the kernel
of reduction to the base ray group is exactly $\prod B_\ell$.
There is no further quotient by the four units.

The map $z\mapsto z/\bar z$ from $B_\ell$ onto $N_\ell$ is surjective.
For split primes this is $(x,y)\mapsto(x/y,y/x)$.
For inert primes it is $z\mapsto z^{1-\ell}$ in the cyclic group
$\mathbb F_{\ell^2}^\times$.
These maps are commutators with complex conjugation, so $N_L\subseteq G_L'$.
The norm on each new residue factor gives the cyclotomic action
$\zeta_\ell\mapsto\zeta_\ell^{\mathrm N(z)}$.
This is also seen by applying arithmetic Frobenius to $\zeta_\ell$;
it is compatible with the ray map by reciprocity.
After quotienting by $N_L$, conjugation acts trivially on the new norm
factors, so
\[
 G_L/N_L\simeq G_0\times\prod_{\ell\mid T_L}\mathbb F_\ell^\times.
\]
Equivalently, the fixed field is $K_0\mathbb Q(\zeta_{T_L})$.
The two factors are disjoint: $K_0$ is unramified outside $2H_0$, while
a nontrivial subextension of the cyclotomic product has nontrivial inertia
at some prime dividing $T_L$. Taking commutators proves both claims about
$G_L'$ and the maximal abelian subfield.

Fix a nonzero norm $s\in\mathbb F_\ell^\times$.
For $z=a+bi$ in that norm fibre,
$\mathrm N(z\pm1)=s+1\pm2a$.
For each sign, vanishing fixes $a$ and leaves at most two values of $b$.
Thus at most four of the $\ell-\chi(\ell)$ elements are excluded.
The conditions are invariant under conjugation, so $C_L$ is a union of
rational conjugacy classes.
In every compatible abelian fibre its numerator in \eqref{eq:v5-mu} is
at least $|C_0|\prod(\ell-\chi(\ell)-4)$, and its denominator is
$|G_0'|\prod(\ell-\chi(\ell))$.
Their quotient proves \eqref{eq:v5-ray-density}.
\end{proof}

\begin{proposition}[Finite-cutoff lower bound]\label{prop:v5-cutoff-lower}
There is a constant $c>0$ independent of $L$ such that the singular series
in Proposition~\ref{prop:v5-artin-patterns} for the fixed field $K_L$ and
the set $C_L$ at every position satisfies
\[
 \mathfrak S_L\geq c(\log L)^{-4k}.
\]
\end{proposition}
\begin{proof}
At the base modulus retain just $u=b_0$, $w=0\pmod{D_0}$.
At each new prime retain all pairs on which the $k$ norm values are units.
Apply \eqref{eq:v5-ray-density} separately to each position.
The norm-fibre denominators cancel as shown in the previous lemma; a
factor inverse to the degree of $K_L$ is not left in the bound.
The normalized proportion of the retained pairs at each new prime is
$\beta_\ell=1+O_k(\ell^{-2})$. Each is positive, since $w=0$, $u\ne0$
is allowed. Their product, together with the tail Euler product, is
bounded below by a fixed positive number. Williams' theorem with modulus
$1$, or its product for the two reduced classes modulo $4$, gives
\[
 \prod_{\ell\mid T_L}(1-4/(\ell-1))^k
       \gg(\log L)^{-4k}.
\]
The base modulus contributes another fixed positive factor.
\end{proof}

\subsection{The final order of limits}\label{subsec:v5-diagonal}

\begin{proof}[Proof of Theorem~\ref{thm:v5-simultaneous-primes}]
Fix the coefficients, $M$, and the base modulus first. Choose one fixed
$L$ so large that
\[
 \frac C{L\log L}<\frac c{32(\log L)^{4k}},
\]
where $C,c$ are from the obstruction and lower-bound propositions.
Now fix $K_L,C_L$. Only after this choice do we use the asymptotic
polynomial theorem or let its small-prime cutoff tend to infinity.

Choose numbers $\tau_j\downarrow0$, all at most $\tau_0$.
Choose successively $N_j$ so large that for $N\geq N_j$
$\omega(N)\leq\tau_j/2$,
that $\tau_j\sqrt N\geq2$, and that the error term in
\eqref{eq:v5-obstruction-tail} is less than
$c\tau_jN^{3/2}/(32(\log L)^{4k})$.
This is possible for each fixed $\tau_j$.
For $N_j\leq N<N_{j+1}$ use $\tau(N)=\tau_j$ and define
$T=\lfloor\tau(N)\sqrt N\rfloor$.
Then $T=o(\sqrt N)$ and $T\geq\tau\sqrt N/2$.
Apply Proposition~\ref{prop:v5-artin-patterns} to this piecewise choice
of $T$. For all sufficiently large $N$ its lower count is at least
$cNT/(4(\log L)^{4k})$ and exceeds the full obstruction tail.
Hence a pattern remains with no forbidden odd
common divisor greater than $L$.
The ray conditions exclude every odd divisor at most $L$ from each
individual adjacent norm. At $7$ the adjacent norms of $1+3i$ are $2$
and $6$; at every other odd base prime the adjacent norms of $i$ are $2$.
Thus no forbidden odd common divisor remains at all.

The residue modulo $16$ gives, on writing $a=9+16A$ and $b=16B$,
\[
 (a+1)^2+b^2\equiv4\pmod8,\qquad
 (a-1)^2+b^2\equiv64\pmod{128}.
\]
These yield the exact two-adic valuations and hence the cross-sign gcd
$4$. The residue modulo $M$ gives coprimality with $M$.
Finally $u\geq N/2$ and $w\leq\tau\sqrt N$ imply
$u/w^2\geq1/(2\tau^2)\to\infty$.
For sufficiently small $\tau_0$, all $u+c_jw^2$ are between fixed positive
multiples of $N$, so they exceed $Y$ when $N$ is large.
The dilation at the start imposes $q\mid w$.
\end{proof}
